\clearpage
\begin{trapbox}
\textbf{\textcolor{CTrap}{易错一：焦点换到纵轴，仍机械代入横坐标}}
\begin{center}
\begin{tikzpicture}[x=1cm,y=1cm,line cap=round,line join=round,font=\small]
\draw[black!25,rounded corners=3pt] (0,0) rectangle (14.4,10);
\node[font=\small\bfseries] at (7.2,9.5) {先看焦点轴：横轴用 $x_0$，纵轴用 $y_0$};
\begin{scope}[shift={(3.6,5.15)},x=0.48cm,y=0.48cm]
\draw[black!35,->] (-5.7,0)--(5.8,0) node[right] {$x$};
\draw[black!35,->] (0,-4.5)--(0,4.6) node[above] {$y$};
\draw[black!65,thick] (0,0) ellipse (5 and 4);
\coordinate (P) at (3,3.2);
\coordinate (Fone) at (-3,0);
\coordinate (Ftwo) at (3,0);
\draw[blue!70!black,line width=1.1pt] (Fone)--(P);
\draw[orange!85!black,line width=1.1pt] (Ftwo)--(P);
\foreach \q in {P,Fone,Ftwo} \fill (\q) circle (1.6pt);
\node[above right] at (P) {$P$};
\node[below left] at (Fone) {$F_1$};
\node[below right] at (Ftwo) {$F_2$};
\end{scope}
\node at (3.6,8.65) {焦点在 $x$ 轴};
\node[blue!70!black] at (3.6,1.8) {$r_1=a+ex_0$};
\node[orange!85!black] at (3.6,1.15) {$r_2=a-ex_0$};
\begin{scope}[shift={(10.8,5.15)},x=0.48cm,y=0.48cm]
\draw[black!35,->] (-4.7,0)--(4.8,0) node[right] {$x$};
\draw[black!35,->] (0,-5.5)--(0,5.6) node[above] {$y$};
\draw[black!65,thick] (0,0) ellipse (4 and 5);
\coordinate (P) at (-3.2,3);
\coordinate (Fone) at (0,-3);
\coordinate (Ftwo) at (0,3);
\draw[blue!70!black,line width=1.1pt] (Fone)--(P);
\draw[orange!85!black,line width=1.1pt] (Ftwo)--(P);
\foreach \q in {P,Fone,Ftwo} \fill (\q) circle (1.6pt);
\node[above left] at (P) {$P$};
\node[below left] at (Fone) {$F_1$};
\node[below right] at (Ftwo) {$F_2$};
\end{scope}
\node at (10.8,8.65) {焦点在 $y$ 轴};
\node[blue!70!black] at (10.8,1.8) {$r_1=a+ey_0$};
\node[orange!85!black] at (10.8,1.15) {$r_2=a-ey_0$};
\node at (7.2,.45) {同一椭圆旋转 $90^\circ$：左右焦点变为下上焦点，距离保持不变。};
\end{tikzpicture}
\end{center}

\textbf{错因与纠正：} 横轴图用 $x_0$，纵轴图用 $y_0$；先定位焦点，再写“左加右减”或“下加上减”。
本图左侧 $P=(3,16/5)$，右侧旋转后 $P'=(-16/5,3)$；两图均有 $r_1=34/5,r_2=16/5$。

\textbf{\textcolor{CTrap}{易错二：加减号写反，或擅自把 $x_0$ 换成 $|x_0|$}}
这样会把左、右两条焦半径对调。用理解页的左右端点检查：
右端点 $P(a,0)$ 应有 $r_1=a+c,r_2=a-c$；左端点正好反过来。
\end{trapbox}
\clearpage
\begin{trapbox}
\textbf{\textcolor{CTrap}{易错三：只检查坐标范围，就对任意点套公式}}
例如本例不在椭圆上的点 $(0,0)$ 虽满足 $x_0\in[-5,5]$，但到两焦点的距离均为 $3$，并非公式给出的 $5$。
\textbf{正确做法：} 必须确认 $P$ 在椭圆上；范围条件只是必要条件。

\medskip
\textbf{\textcolor{CTrap}{易错四：准线距离直接去绝对值，不说明依据}}
右准线距离本为 $|x_0-a^2/c|$。因为 $P$ 在椭圆上且 $a^2/c>a$，
才可写为 $a^2/c-x_0$。准线图中两条准线都在长轴端点外侧。

\medskip
\textbf{\textcolor{CTrap}{易错五：由平方相等直接开方，漏掉正性}}
由 $r_2^2=(a-ex_0)^2$ 先得 $r_2=|a-ex_0|$；
再由 $|x_0|\le a,e<1$ 知 $a\pm ex_0\ge a-c>0$，才能去掉绝对值。

\medskip
\textbf{\textcolor{CTrap}{易错六：仅用三个特殊位置替代一般证明}}
三态图可以检查符号、观察极值，但不能证明全椭圆上的公式。
一般证明仍需准线定义、第一定义加勾股平方差，或代入椭圆方程配方。
\medskip
\textbf{\textcolor{CTrap}{易错七：单变量法中误认为 $t$ 一定大于零}}
设 $|PF_1|=a+t$、$|PF_2|=a-t$ 时，$t$ 是\textbf{带符号的差量}，
不一定是正数。$x_0<0$ 时 $t=ex_0<0$，此时左焦半径反而短；
不能擅自改设 $a+|t|,a-|t|$。真正需要保证的是 $a\pm t>0$。
\end{trapbox}
